\chapter{Short-Term Futures Strategies}\label{ch:s1:short-term-futures-strategies}

Lucca and Moench found large average excess returns on US equities in the hours before scheduled Federal Reserve policy decisions: returns that had
grown over time, accounted for a sizable share of the market's annual return, appeared in other countries' indices, and had no counterpart before other
macroeconomic announcements. Once published, such an effect is everybody's. Short-term futures strategies live on patterns like it, inside the
trading day: the first half hour, the moves that bounce back within minutes, the hours before a scheduled release. On this chapter's synthetic
sessions the pre-announcement hold earns 40.6 basis points per event before its planted crowding and 18.8 after, a Sharpe ratio of 1.21 falling to
0.58; a fade of large five-minute moves earns a Sharpe ratio of 2.86 before costs and $-1.50$ after half a spread each way. The build is
\texttt{firm.futintraday}.

\section{Sessions in bars}

The index future's regular session is simulated in one-minute bars (\cref{lst:s1:short-term-futures-strategies:sessions}): ten years of 390-minute
sessions, 16.2\% volatility a year with a quarter of the variance in the overnight gap, a U-shaped intraday profile (twice as volatile at the open and
close as at midday), a daily trend drift with a standard deviation of 0.4\% over the session, and a bounce of $-0.1$ in one-minute returns. Eight times a
year a scheduled announcement arrives at minute 270; a drift of 0.5\% is planted from the previous day's midday to the announcement, and from the sixth
year it is cut to 30\% of that: the crowding that follows publication. Book~7's message-level \texttt{firm.tape} and event-driven \texttt{firm.evbt} are
the tools for queue-level questions; for ten years of sessions, bars suffice, and the tape supplies the cost. On one simulated tape session (a price of
100.00, a tick of 0.01, so a tick is a basis point) the time-weighted quoted spread is 1.08 ticks; each trade here pays half of it, 0.54 basis points.

\section{Opening ranges and intraday mean reversion}

\begin{definition}[Opening range breakout]\label{def:s1:short-term-futures-strategies:orb}
An \emph{opening range breakout}\index{opening range breakout} records the highest and lowest prices of the first minutes of the session (the opening
range) and takes a position in the direction of the first later move outside it, held to the close or a stop.
\end{definition}

\begin{definition}[Intraday mean reversion]\label{def:s1:short-term-futures-strategies:mr}
\emph{Intraday mean reversion}\index{intraday mean reversion} is the tendency of short intraday price moves to be partly reversed within minutes, from
bid--ask bounce and liquidity providers' inventory; a mean-reversion strategy fades moves larger than a threshold and exits after a fixed holding time.
\end{definition}

The two rules bet on opposite features of the synthetic session (\cref{lst:s1:short-term-futures-strategies:rules}). The breakout bets on the trend day:
once the session has moved outside its first minutes' range, the day's drift is more likely in that direction. The fade bets on the bounce: after the
price has moved more than $z$ standard deviations in five minutes, sell the move and exit five minutes later.

\begin{center}
\small
\begin{tabular}{@{}lccc@{}}
\toprule
opening range (minutes) & 15 & 30 & 60\\
\midrule
breakout: Sharpe ratio before costs & 0.56 & 0.77 & 1.49\\
breakout: Sharpe ratio after costs & 0.35 & 0.55 & 1.25\\
breakout: mean daily return after costs (bp) & 1.88 & 2.77 & 5.62\\
\bottomrule
\end{tabular}
\end{center}

A longer range gives a later, surer read of the day's direction: the planted drift is spread evenly over the session, so an hour reveals more of it
than a quarter hour, and the breakout trades once a day at most, so costs matter little. The fade is the opposite (\cref{fig:s1:short-term-futures-strategies:fade}):
the bounce is strong enough for Sharpe ratios of 0.59 to 3.78 before costs, and at 0.54 basis points each way every threshold loses. The bounce gives
back a tenth of the last minute's move, a small part of a five-minute move, and that part is smaller than a spread.

\begin{omfigure}
\begin{tikzpicture}
\begin{axis}[width=10cm,height=5cm,ybar,bar width=12pt,xlabel={\small threshold (standard deviations of a five-minute move)},
  ylabel={\small Sharpe ratio},tick label style={font=\footnotesize},symbolic x coords={1.5,2.0,2.5,3.0},xtick=data,ymin=-4,ymax=4.5,
  enlarge x limits=0.2,grid=major,grid style={gray!25},legend columns=2,legend style={font=\scriptsize,at={(0.5,-0.3)},anchor=north,
  draw=none,fill=none,/tikz/every even column/.append style={column sep=8pt}},area legend,table/col sep=comma]
\addplot[fill=omExo!50,draw=omExo] table[x=z,y=gross]{figdata/strategies-1/25-short-term-futures-strategies/fade.csv};
\addplot[fill=omThm!60,draw=omThm] table[x=z,y=net]{figdata/strategies-1/25-short-term-futures-strategies/fade.csv};
\legend{before costs,after half a spread each way}
\end{axis}
\end{tikzpicture}

\omcaption{The intraday fade on ten years of synthetic one-minute sessions: Sharpe ratios before and after costs of fading five-minute moves beyond a
threshold and holding five minutes. The planted bounce is real; half the quoted spread of Book~7's tape, each way, is larger. Data:
\texttt{s1\_futures.fade\_table}.}\label{fig:s1:short-term-futures-strategies:fade}
\end{omfigure}

\section{Scheduled announcements}

\begin{definition}[Pre-announcement drift]\label{def:s1:short-term-futures-strategies:pre}
A \emph{pre-announcement drift}\index{pre-announcement drift} is an average price move in the hours before a scheduled announcement, in a direction
that does not depend on the announcement's content; the pre-FOMC drift in US equities is the best-known case.
\end{definition}

\begin{dated}{2026-09}{FOMC meetings in 2026}\label{dat:s1:short-term-futures-strategies:fomc}
The Federal Reserve's calendar lists eight FOMC meetings in 2026: January 27--28, March 17--18, April 28--29, June 16--17, July 28--29, September
15--16, October 27--28 and December 8--9; the March, June, September and December meetings include a Summary of Economic Projections.
\end{dated}

Lucca and Moench's drift is unusual in two ways: it precedes the news instead of following it, and it has a sign regardless of what is announced. It
also sits in equities only; they found none in Treasuries or money-market futures, and none before other major announcements. That makes it a risk
premium or a behavioural pattern, not information, and it makes it easy to trade: hold the equity index future from the day before to the
announcement.

\begin{center}
\small
\begin{tabular}{@{}lcc@{}}
\toprule
pre-announcement hold, after costs & years 1--5 & years 6--10 (crowded)\\
\midrule
events & 39 & 40\\
mean return per event (bp) & 40.6 & 18.8\\
t-statistic & 2.7 & 1.3\\
Sharpe ratio a year (eight events) & 1.21 & 0.58\\
\bottomrule
\end{tabular}
\end{center}

The planted drift is 50 basis points before the crowding and 15 after. The measured means, 40.6 and 18.8, are within a standard error of them: with a
day's volatility of about 1\% per event, 39 events measure the mean to about 16 basis points (\cref{fig:s1:short-term-futures-strategies:path}). That is
the practical difficulty with event strategies: eight events a year is a small sample, and a decay to a third of the effect is barely visible in five
years. The published drift is a strong prior, and a live book can only confirm it slowly.

\begin{omfigure}
\begin{tikzpicture}
\begin{axis}[width=11cm,height=5.2cm,xlabel={\small minutes from the previous day's midday},ylabel={\small mean cumulative return (bp)},
  tick label style={font=\footnotesize},xmin=0,xmax=590,ymin=-20,ymax=70,grid=major,grid style={gray!25},legend pos=north west,
  legend style={font=\scriptsize},table/col sep=comma]
\addplot[omThm,thick] table[x=minute,y=before]{figdata/strategies-1/25-short-term-futures-strategies/announcement.csv};
\addplot[omDef,thick] table[x=minute,y=after]{figdata/strategies-1/25-short-term-futures-strategies/announcement.csv};
\addplot[black!50,dashed,forget plot] coordinates {(466,-20) (466,70)};
\legend{years 1--5,years 6--10 (crowded)}
\end{axis}
\end{tikzpicture}

\omcaption{The average path of the synthetic index future around scheduled announcements, from the previous day's midday (the overnight gap counted as
one step) to the announcement day's close; the dashed line marks the announcement. Planted: 50 basis points before it in the first five years, 15 in
the last five. Data: \texttt{s1\_futures.announcement\_path}.}\label{fig:s1:short-term-futures-strategies:path}
\end{omfigure}

\section{What makes a short-term edge survive}

Three features decide. Turnover against edge: the fade trades about seven times a day for a fraction of a basis point each and loses; the breakout trades
once and keeps most of its edge. Sample size against decay: an event strategy with eight events a year cannot see its own crowding for years. And the
mechanism: a bounce is real but belongs to the liquidity providers who earn the spread; a trend day's direction and a scheduled drift are harder to
explain and therefore harder to be sure of. The overnight gap and the reaction after announcements are further candidates; the strategy files record
them without a verified performance figure.

\section{Strategy files}

\begin{strategyfile}{Opening range breakout}\label{strat:s1:short-term-futures-strategies:orb}
\sfield{Who pays you, and why} Traders who react late to the day's direction.
\sfield{Instruments and venues} Liquid index, rates and commodity futures.
\sfield{Signal} The first move outside the opening range.
\sfield{Sizing and execution} One position a day, held to the close or a stop; volatility-scaled.
\sfield{Costs} Low: one round trip a day.
\sfield{How it dies} Days without trends; crowding at the range's edges.
\sfield{Horizon, capacity, infrastructure} Hours; intraday bars.
\sfield{Backtest honestly} Fills at the breakout level plus slippage; the range's length chosen before the test.
\sfield{Sources} No performance figure verified; this chapter: 0.55 after costs with a 30-minute range, synthetic.
\end{strategyfile}

\begin{strategyfile}{Intraday futures mean reversion}\label{strat:s1:short-term-futures-strategies:mr}
\sfield{Who pays you, and why} Impatient traders who push prices briefly; in practice the liquidity providers collect this.
\sfield{Instruments and venues} Liquid futures.
\sfield{Signal} Five-minute moves beyond a threshold.
\sfield{Sizing and execution} Fade, hold minutes; passive entry if possible.
\sfield{Costs} Decisive: aggressive orders pay the spread on every trade.
\sfield{How it dies} Costs; faster liquidity providers.
\sfield{Horizon, capacity, infrastructure} Minutes; queue-aware execution (Book~7's level-2 backtester).
\sfield{Backtest honestly} The spread on every aggressive fill; queue position for passive ones.
\sfield{Sources} This chapter: 2.86 before costs, $-1.50$ after, synthetic.
\end{strategyfile}

\begin{strategyfile}{Pre-FOMC drift}\label{strat:s1:short-term-futures-strategies:prefomc}
\sfield{Who pays you, and why} Unclear: a premium for bearing announcement risk, or investors' behaviour before policy news.
\sfield{Instruments and venues} Equity index futures.
\sfield{Signal} The FOMC calendar.
\sfield{Sizing and execution} Long from the day before to the announcement.
\sfield{Costs} One round trip per meeting.
\sfield{How it dies} Publication and crowding.
\sfield{Horizon, capacity, infrastructure} A day, eight times a year; large capacity.
\sfield{Backtest honestly} Meeting dates and statement times as scheduled; unscheduled meetings excluded; enough events to judge decay.
\sfield{Sources} Lucca and Moench (2015); this chapter: 40.6 basis points per event before crowding, 18.8 after.
\end{strategyfile}

\begin{strategyfile}{Post-announcement momentum}\label{strat:s1:short-term-futures-strategies:post}
\sfield{Who pays you, and why} Traders who adjust to scheduled news gradually.
\sfield{Instruments and venues} Rates and index futures around data releases.
\sfield{Signal} The direction of the first minutes after a release, relative to the surprise.
\sfield{Sizing and execution} Follow, hold hours.
\sfield{Costs} Wide spreads at releases.
\sfield{How it dies} Machines that finish the reaction in seconds (chapter~17).
\sfield{Horizon, capacity, infrastructure} Minutes to hours; low-latency release data.
\sfield{Backtest honestly} Release timestamps to the second; spreads at the release.
\sfield{Sources} No performance figure verified.
\end{strategyfile}

\begin{strategyfile}{Overnight gap fade}\label{strat:s1:short-term-futures-strategies:gap}
\sfield{Who pays you, and why} Overnight orders that overshoot at the open.
\sfield{Instruments and venues} Index futures at the cash open.
\sfield{Signal} The gap between the previous close and the open, against its volatility.
\sfield{Sizing and execution} Fade large gaps, exit within the morning.
\sfield{Costs} Wide opening spreads.
\sfield{How it dies} Gaps that carry news continue.
\sfield{Horizon, capacity, infrastructure} Hours.
\sfield{Backtest honestly} The tradeable opening price, not the print; news days separated.
\sfield{Sources} No performance figure verified; chapter~13's \omterm{def:s1:intraday-patterns:returns}{overnight return}.
\end{strategyfile}

\section{Tutorial: before the announcement}\label{tut:s1:short-term-futures-strategies:tut}

\begin{tutorial}
\textbf{Goal.} Simulate ten years of one-minute futures sessions with a trend-day drift, a bounce and a crowded \omterm{def:s1:short-term-futures-strategies:pre}{pre-announcement drift}; take the cost of
a trade from Book~7's tape; test the breakout, the fade and the pre-announcement hold. \textbf{End state:} the tables and the two figures.
\begin{tutsteps}
\item \textbf{Sessions}: the U-shaped profile, the bounce, the trend day and the planted announcements.
\omcode{code/firm/futintraday/firm_futintraday.py}{50}{71}{Ten years of one-minute sessions.}\label{lst:s1:short-term-futures-strategies:sessions}
\item \textbf{The rules}: the \omterm{def:s1:short-term-futures-strategies:orb}{opening range breakout} and the fade.
\omcode{code/firm/futintraday/firm_futintraday.py}{74}{107}{Opening range breakout and intraday fade.}\label{lst:s1:short-term-futures-strategies:rules}
\item \textbf{Run} \texttt{tape\_cost()}, \texttt{orb\_table()}, \texttt{fade\_table()}, \texttt{announcements()} and \texttt{announcement\_path()}, and
\texttt{fig\_futures.py}.
\end{tutsteps}
\textbf{What to change next.} Enter the fade passively at the far side of the spread and model the fill probability; run the pre-announcement strategy on
forty years of events to see how long the crowding takes to detect; replace the bars with firm.tape sessions for the announcement minutes.
\end{tutorial}

\section{Build: intraday futures toolkit}\label{bld:s1:short-term-futures-strategies:futintraday}

\begin{build}
\bfield{Purpose} Bar-level futures sessions with planted features, and the breakout, fade and pre-announcement strategies with costs.
\bfield{Interface} \texttt{SessionConfig(\dots)}, \texttt{simulate\_sessions(cfg, rng)}, \texttt{orb(S, window, cost)}, \texttt{fade(S, lookback, z, hold,
cost)}, \texttt{pre\_announcement(S, cost)}.
\bfield{Rules} Decisions from bars already closed; one position at a time; two costs per round trip.
\bfield{Acceptance tests} \texttt{code/firm/futintraday/tests/}: the session's volatility and bounce, and a breakout and a pre-announcement hold by hand.
\bfield{Stretch} Passive fills; gap strategies; announcements with surprises and post-announcement dynamics.
\end{build}

\begin{omsources}
\item D.~O.~Lucca and E.~Moench, ``The pre-FOMC announcement drift'', \emph{Journal of Finance} 70(1), 2015.
\item Board of Governors of the Federal Reserve System, FOMC meeting calendars.
\end{omsources}

\section{Exercises}

\begin{exercise}[$\star$]\label{exo:s1:short-term-futures-strategies:1}
At a price of 100.00 and a tick of 0.01, how many basis points is a tick? What is the round-trip cost at a quoted spread of 1.08 ticks?
\end{exercise}

\begin{exercise}[$\star$]\label{exo:s1:short-term-futures-strategies:2}
The planted drift of 50 basis points is cut to 30\%. What is left, and what annual Sharpe ratio multiplier does trading eight events a year give a
per-event Sharpe ratio?
\end{exercise}

\begin{exercise}[$\star$]\label{exo:s1:short-term-futures-strategies:3}
Why does a longer opening range work better here?
\end{exercise}

\begin{exercise}[$\star\star$]\label{exo:s1:short-term-futures-strategies:4}
With a per-event standard deviation of about 1\%, what is the standard error of a mean over 39 events, and what t-statistic does a true drift of 50
basis points have on average?
\end{exercise}

\begin{exercise}[$\star\star$]\label{exo:s1:short-term-futures-strategies:5}
Why does the fade lose after costs at every threshold, even though it wins before?
\end{exercise}

\begin{exercise}[$\star\star$]\label{exo:s1:short-term-futures-strategies:6}
A bounce of $-0.1$ in one-minute returns reduces a day's session variance. By how much does it reduce a day's volatility when the session carries three
quarters of the variance?
\end{exercise}

\begin{exercise}[$\star\star\star$]\label{exo:s1:short-term-futures-strategies:7}
\emph{Coding.} Run \texttt{fade\_stats(2.0)}. How much does the fade earn a day before costs, how many trades does it make, and what cost per side
would it break even at?
\end{exercise}

\begin{exercise}[$\star\star\star$]\label{exo:s1:short-term-futures-strategies:8}
\emph{Find the flaw.} ``The pre-FOMC drift earned 40 basis points per meeting over the last five years, so the edge is intact.''
\end{exercise}

\section{Problem: Before the Announcement}

\begin{problem}[{Weekend problem --- three intraday edges}]\label{pb:s1:short-term-futures-strategies:1}
The chapter's synthetic sessions and the public record.

\medskip\noindent\textbf{Part I --- Sessions.}
\begin{enumerate}
\item Describe the bar model and what is planted in it.
\item How is the cost of a trade obtained from Book~7's tape?
\item Define the \omterm{def:s1:short-term-futures-strategies:orb}{opening range breakout} and \omterm{def:s1:short-term-futures-strategies:mr}{intraday mean reversion}.
\item What feature of the session does each bet on?
\end{enumerate}

\medskip\noindent\textbf{Part II --- Rules.}
\begin{enumerate}[resume]
\item Give the breakout's Sharpe ratios by range length.
\item Give the fade's Sharpe ratios before and after costs.
\item Why does one survive costs and not the other?
\item Who earns the bounce in real markets?
\end{enumerate}

\medskip\noindent\textbf{Part III --- Announcements.}
\begin{enumerate}[resume]
\item Define the \omterm{def:s1:short-term-futures-strategies:pre}{pre-announcement drift}; what did Lucca and Moench find?
\item What does the dated box list?
\item Give the per-event returns and Sharpe ratios before and after crowding.
\item Why is crowding hard to detect?
\end{enumerate}

\medskip\noindent\textbf{Part IV --- The verdict.}
\begin{enumerate}[resume]
\item State the \emph{named result}: the \omterm{def:s1:short-term-futures-strategies:pre}{pre-announcement drift}'s Sharpe ratio before and after the planted crowding.
\item List the three features that decide whether a short-term edge survives.
\item How would you backtest the pre-FOMC drift honestly?
\item Which strategy file depends most on latency?
\item How does this chapter relate to chapter~13's intraday patterns?
\item What would you need to run the fade profitably?
\item Why is the pre-FOMC drift not information?
\item In one sentence: what do short-term futures strategies sell?
\end{enumerate}
\end{problem}

\section{Interview questions}

\begin{interviewq}[$\star$ \iqroles{trader}]\label{iq:s1:short-term-futures-strategies:1}
What is an \omterm{def:s1:short-term-futures-strategies:orb}{opening range breakout}, and when does it fail?
\end{interviewq}

\begin{interviewq}[$\star\star$ \iqroles{researcher}]\label{iq:s1:short-term-futures-strategies:2}
How would you test whether a published \omterm{def:s1:short-term-futures-strategies:pre}{pre-announcement drift} still exists?
\end{interviewq}

\begin{interviewq}[$\star\star$ \iqroles{researcher}]\label{iq:s1:short-term-futures-strategies:3}
Your intraday strategy has a Sharpe ratio of 3 before costs. What do you check first?
\end{interviewq}

\begin{interviewq}[$\star\star$ \iqroles{developer}]\label{iq:s1:short-term-futures-strategies:4}
What does a backtester need to simulate an intraday futures strategy faithfully?
\end{interviewq}

\begin{interviewq}[$\star\star$ \iqroles{risk}]\label{iq:s1:short-term-futures-strategies:5}
What are the risks of holding index futures into scheduled announcements?
\end{interviewq}

\begin{interviewq}[$\star\star\star$ \iqroles{researcher}]\label{iq:s1:short-term-futures-strategies:6}
One-minute returns are $r_t = e_t + \theta e_{t-1}$ with independent $e_t$ of variance $\sigma^2$. A trader fades the last minute's return and holds one
minute. Derive the expected P\&L per trade and the cost per trade at which it breaks even.
\end{interviewq}
